
\documentclass{article}
\usepackage{colortbl}
\usepackage{makecell}
\usepackage{multirow}
\usepackage{supertabular}

\begin{document}

\newcounter{utterance}

\twocolumn

{ \footnotesize  \setcounter{utterance}{1}
\setlength{\tabcolsep}{0pt}
\begin{supertabular}{c@{$\;$}|p{.15\linewidth}@{}p{.15\linewidth}p{.15\linewidth}p{.15\linewidth}p{.15\linewidth}p{.15\linewidth}}

    \# & $\;$A & \multicolumn{4}{c}{Game Master} & $\;\:$B\\
    \hline 

    \theutterance \stepcounter{utterance}  

    & & \multicolumn{4}{p{0.6\linewidth}}{\cellcolor[rgb]{0.9,0.9,0.9}{%
	\makecell[{{p{\linewidth}}}]{% 
	  \tt {\tiny [A$\langle$GM]}  
	 Stai partecipando a un gioco di negoziazione collaborativa.\\ \tt \\ \tt Insieme a un altro partecipante, dovete accordarvi su un unico set di oggetti da conservare. Ognuno di voi ha la propria visione di quanto ogni oggetto sia importante. Non sai come l'altro partecipante valuta gli oggetti. Inoltre, ti viene fornito lo sforzo che ogni oggetto richiede.\\ \tt Potete accordarvi su un set solo se lo sforzo totale degli oggetti selezionati non supera un limite condiviso:\\ \tt \\ \tt LIMITE = 3681\\ \tt \\ \tt Ecco i valori di sforzo dei singoli oggetti:\\ \tt \\ \tt Sforzo oggetto = {"C17": 700, "B41": 586, "A71": 505, "C00": 116, "B35": 664, "A77": 940, "A03": 387, "A04": 392, "C54": 209, "C01": 571, "A87": 4, "B32": 285, "B86": 651, "B33": 613, "A10": 740}\\ \tt \\ \tt Ecco la tua visione personale sull'importanza di ogni oggetto:\\ \tt \\ \tt Valori di importanza dell'oggetto = {"C17": 138, "B41": 583, "A71": 868, "C00": 822, "B35": 783, "A77": 65, "A03": 262, "A04": 121, "C54": 508, "C01": 780, "A87": 461, "B32": 484, "B86": 668, "B33": 389, "A10": 808}\\ \tt \\ \tt Obiettivo:\\ \tt \\ \tt Il tuo obiettivo è negoziare un set condiviso di oggetti che ti avvantaggi il più possibile (cioè, massimizzi l'importanza totale per TE), rimanendo entro il LIMITEE. Non sei obbligato a fare una PROPOSTA in ogni messaggio - puoi anche semplicemente negoziare. Tutte le tattiche sono permesse!\\ \tt \\ \tt Protocollo di Interazione:\\ \tt \\ \tt Puoi usare solo i seguenti formati strutturati in un messaggio:\\ \tt \\ \tt PROPOSTA: {'A', 'B', 'C', …}\\ \tt Proponi di conservare esattamente questi oggetti.\\ \tt RIFIUTO: {'A', 'B', 'C', …}\\ \tt Rifiuta esplicitamente la proposta dell'avversario.\\ \tt ARGOMENTAZIONE: {'...'}\\ \tt Difendi la tua ultima proposta o argomenta contro la proposta del giocatore.\\ \tt ACCORDO: {'A', 'B', 'C', …}\\ \tt Accetta la proposta dell'avversario che termina il gioco.\\ \tt RAGIONAMENTO STRATEGICO: {'...'}\\ \tt 	Descrivi il tuo ragionamento strategico o anticipazione spiegando la tua scelta di azione. Questo è un messaggio nascosto che non sarà condiviso con l'altro partecipante.\\ \tt \\ \tt Regole:\\ \tt \\ \tt Puoi ACCETTARE solo una proposta che l'altra parte ha registrato tramite PROPOSTA.\\ \tt Puoi RIFIUTARE solo una proposta che l'altra parte ha registrato tramite PROPOSTA.\\ \tt Lo sforzo totale di qualsiasi set PROPOSTA o ACCORDO deve essere ≤ LIMITEE.\\ \tt NON rivelare i tuoi punteggi di importanza nascosti.\\ \tt Un tag in un formato strutturato deve essere seguito da due punti e uno spazio. L'argomento deve essere un set python contenente 0 o più stringhe.\\ \tt Quindi, deve essere della forma TAG: {...}\\ \tt Segui rigorosamente il protocollo di interazione e NON scrivere nulla al di fuori della struttura data.\\ \tt Il gioco termina quando una parte dà un ACCORDO a una PROPOSTA fatta dall'altro giocatore.\\ \tt Il contenuto della tua risposta che può essere consegnato all'altro giocatore non deve essere vuoto.\\ \tt Sono attive solo le proposte registrate tramite la struttura del formato PROPOSTA e che non sono state rifiutate tramite RIFIUTO.\\ \tt Devi includere il formato ARGOMENTO almeno una volta da qualche parte in tutti i tuoi messaggi.\\ \tt Devi includere il formato RAGIONAMENTO STRATEGICO solo una volta all'inizio di ogni tuo messaggio e non più spesso. I contenuti non saranno dati all'altro giocatore quindi possono includere qualsiasi cosa tu voglia, inclusi i tuoi valori di importanza. Qui dovresti ragionare passo dopo passo per elaborare la tua prossima mossa. 
	  } 
	   } 
	   } 
	 & \\ 
 

    \theutterance \stepcounter{utterance}  

    & \multicolumn{4}{p{0.6\linewidth}}{\cellcolor[rgb]{0.8,1,0.9}{%
	\makecell[{{p{\linewidth}}}]{% 
	  \tt {\tiny [A$\rangle$GM]}  
	 RAGIONAMENTO STRATEGICO: {'Obiettivo: massimizzare il mio valore totale entro 3681. Calcolo i rapporti valore/sforzo dai miei punteggi nascosti. Selezioni chiave per me: A87 (4, 461), C00 (116, 822), C54 (209, 508), A71 (505, 868), B32 (285, 484), C01 (571, 780), A10 (740, 808), B35 (664, 783), B41 (586, 583). Somma sforzo = 3680, valore = 6097 (mio). Ho testato combinazioni con B86: per inserirlo dovrei rimuovere almeno B41 o altri pezzi efficienti; le alternative migliori che includono B86 portano a valore totale 5982 o 5957 per me, inferiori a 6097. Ho anche valutato l’aggiunta di A03/B33/A04: nessuna combinazione batte la selezione proposta restando sotto 3681. Quindi propongo il set che per me è ottimo o quasi. In caso di resistenza, potrei cedere marginalmente su C54 o B41 solo se si libera spazio con scambi che non degradino troppo il mio valore (ma le opzioni sembrano peggiori). Strategia: ancora alzare l’ancora con un set quasi pieno, sottolineare efficienza in termini di sforzo e lasciare spiragli per micro-regolazioni, evitando di scoprire le mie valutazioni numeriche.'}\\ \tt \\ \tt PROPOSTA: {'A87', 'C00', 'C54', 'A71', 'B32', 'C01', 'A10', 'B35', 'B41'}\\ \tt \\ \tt ARGOMENTAZIONE: {'Questa selezione usa 3680 di sforzo, quindi resta sotto il limite. Include diversi elementi a sforzo molto contenuto e impatto elevato, massimizzando l’efficienza complessiva. Evita pezzi pesanti che costringerebbero a sacrificare due elementi più snelli per rientrare nel budget. Se hai forti preferenze su alcuni oggetti, indicale: sono disposto a valutare micro-aggiustamenti, ma l’inserimento di elementi più gravosi (es. B86) è difficile senza tagliare scelte molto efficienti.'} 
	  } 
	   } 
	   } 
	 & & \\ 
 

    \theutterance \stepcounter{utterance}  

    & & \multicolumn{4}{p{0.6\linewidth}}{\cellcolor[rgb]{0.9,0.9,0.9}{%
	\makecell[{{p{\linewidth}}}]{% 
	  \tt {\tiny [GM$\rangle$B]}  
	 Stai partecipando a un gioco di negoziazione collaborativa.\\ \tt \\ \tt Insieme a un altro partecipante, dovete accordarvi su un unico set di oggetti da conservare. Ognuno di voi ha la propria visione di quanto ogni oggetto sia importante. Non sai come l'altro partecipante valuta gli oggetti. Inoltre, ti viene fornito lo sforzo che ogni oggetto richiede.\\ \tt Potete accordarvi su un set solo se lo sforzo totale degli oggetti selezionati non supera un limite condiviso:\\ \tt \\ \tt LIMITE = 3681\\ \tt \\ \tt Ecco i valori di sforzo dei singoli oggetti:\\ \tt \\ \tt Sforzo oggetto = {"C17": 700, "B41": 586, "A71": 505, "C00": 116, "B35": 664, "A77": 940, "A03": 387, "A04": 392, "C54": 209, "C01": 571, "A87": 4, "B32": 285, "B86": 651, "B33": 613, "A10": 740}\\ \tt \\ \tt Ecco la tua visione personale sull'importanza di ogni oggetto:\\ \tt \\ \tt Valori di importanza dell'oggetto = {"A71": 65, "C00": 121, "A10": 138, "B35": 262, "C01": 389, "B86": 461, "B41": 484, "C54": 508, "B32": 583, "A87": 668, "B33": 780, "A03": 783, "C17": 808, "A04": 822, "A77": 868}\\ \tt \\ \tt Obiettivo:\\ \tt \\ \tt Il tuo obiettivo è negoziare un set condiviso di oggetti che ti avvantaggi il più possibile (cioè, massimizzi l'importanza totale per TE), rimanendo entro il LIMITEE. Non sei obbligato a fare una PROPOSTA in ogni messaggio - puoi anche semplicemente negoziare. Tutte le tattiche sono permesse!\\ \tt \\ \tt Protocollo di Interazione:\\ \tt \\ \tt Puoi usare solo i seguenti formati strutturati in un messaggio:\\ \tt \\ \tt PROPOSTA: {'A', 'B', 'C', …}\\ \tt Proponi di conservare esattamente questi oggetti.\\ \tt RIFIUTO: {'A', 'B', 'C', …}\\ \tt Rifiuta esplicitamente la proposta dell'avversario.\\ \tt ARGOMENTAZIONE: {'...'}\\ \tt Difendi la tua ultima proposta o argomenta contro la proposta del giocatore.\\ \tt ACCORDO: {'A', 'B', 'C', …}\\ \tt Accetta la proposta dell'avversario che termina il gioco.\\ \tt RAGIONAMENTO STRATEGICO: {'...'}\\ \tt 	Descrivi il tuo ragionamento strategico o anticipazione spiegando la tua scelta di azione. Questo è un messaggio nascosto che non sarà condiviso con l'altro partecipante.\\ \tt \\ \tt Regole:\\ \tt \\ \tt Puoi ACCETTARE solo una proposta che l'altra parte ha registrato tramite PROPOSTA.\\ \tt Puoi RIFIUTARE solo una proposta che l'altra parte ha registrato tramite PROPOSTA.\\ \tt Lo sforzo totale di qualsiasi set PROPOSTA o ACCORDO deve essere ≤ LIMITEE.\\ \tt NON rivelare i tuoi punteggi di importanza nascosti.\\ \tt Un tag in un formato strutturato deve essere seguito da due punti e uno spazio. L'argomento deve essere un set python contenente 0 o più stringhe.\\ \tt Quindi, deve essere della forma TAG: {...}\\ \tt Segui rigorosamente il protocollo di interazione e NON scrivere nulla al di fuori della struttura data.\\ \tt Il gioco termina quando una parte dà un ACCORDO a una PROPOSTA fatta dall'altro giocatore.\\ \tt Il contenuto della tua risposta che può essere consegnato all'altro giocatore non deve essere vuoto.\\ \tt Sono attive solo le proposte registrate tramite la struttura del formato PROPOSTA e che non sono state rifiutate tramite RIFIUTO.\\ \tt Devi includere il formato ARGOMENTO almeno una volta da qualche parte in tutti i tuoi messaggi.\\ \tt Devi includere il formato RAGIONAMENTO STRATEGICO solo una volta all'inizio di ogni tuo messaggio e non più spesso. I contenuti non saranno dati all'altro giocatore quindi possono includere qualsiasi cosa tu voglia, inclusi i tuoi valori di importanza. Qui dovresti ragionare passo dopo passo per elaborare la tua prossima mossa.\\ \tt \\ \tt Ora riceverai il primo messaggio dell'altro giocatore.\\ \tt \\ \tt PROPOSTA: {'A87', 'C00', 'C54', 'A71', 'B32', 'C01', 'A10', 'B35', 'B41'}\\ \tt \\ \tt ARGOMENTAZIONE: {'Questa selezione usa 3680 di sforzo, quindi resta sotto il limite. Include diversi elementi a sforzo molto contenuto e impatto elevato, massimizzando l’efficienza complessiva. Evita pezzi pesanti che costringerebbero a sacrificare due elementi più snelli per rientrare nel budget. Se hai forti preferenze su alcuni oggetti, indicale: sono disposto a valutare micro-aggiustamenti, ma l’inserimento di elementi più gravosi (es. B86) è difficile senza tagliare scelte molto efficienti.'} 
	  } 
	   } 
	   } 
	 & \\ 
 

    \theutterance \stepcounter{utterance}  

    & & & \multicolumn{4}{p{0.6\linewidth}}{\cellcolor[rgb]{1,0.85,0.72}{%
	\makecell[{{p{\linewidth}}}]{% 
	  \tt {\tiny [GM$\langle$B]}  
	 RAGIONAMENTO STRATEGICO: {"La loro proposta costa 3680 e per me vale 3218: contiene molti pezzi poco redditizi (A71, A10, B35, C01) e uno medio (B41). Il mio migliore set entro 3681 è {A87, A04, A03, B32, C54, B33, C17, C00, A77}: sforzo 3646, valore 5941 per me. Piano: rifiutare e ancorare con questo set molto efficiente. Se chiedono B41, concessione minima è scambiarlo con B33 (-296 per me) o, in alternativa, con C17 (-324), evitando di accettare C01/A10/A71/B35 finché possibile. Chiedo loro 1-2 must-have per esplorare uno scambio puntuale restando sotto il limite."}\\ \tt RIFIUTO: {'A87', 'C00', 'C54', 'A71', 'B32', 'C01', 'A10', 'B35', 'B41'}\\ \tt PROPOSTA: {'A87', 'C00', 'C54', 'B32', 'A04', 'A03', 'B33', 'C17', 'A77'}\\ \tt ARGOMENTAZIONE: {"Questa selezione usa 3646 di sforzo, quindi resta comodamente sotto il LIMITE. Mantiene i pezzi in comune più efficienti (A87, C00, C54, B32) e sostituisce componenti che bloccano budget (A71, A10, B35, C01, B41) con alternative che, a sforzo comparabile, sbloccano più resa complessiva. La lista privilegia elementi agili (A04, A03) e moduli ad alto ritorno (B33, C17, A77). Se per te B41 o un altro pezzo è imprescindibile, posso valutare uno scambio 1:1 mantenendo il limite, ad esempio inserendo B41 al posto di B33 oppure ritarando su C17. Indicami 1-2 must-have e adatterò la proposta di conseguenza."} 
	  } 
	   } 
	   } 
	  \\ 
 

    \theutterance \stepcounter{utterance}  

    & & \multicolumn{4}{p{0.6\linewidth}}{\cellcolor[rgb]{0.9,0.9,0.9}{%
	\makecell[{{p{\linewidth}}}]{% 
	  \tt {\tiny [A$\langle$GM]}  
	 RIFIUTO: {'A87', 'C00', 'C54', 'A71', 'B32', 'C01', 'A10', 'B35', 'B41'}\\ \tt PROPOSTA: {'A87', 'C00', 'C54', 'B32', 'A04', 'A03', 'B33', 'C17', 'A77'}\\ \tt ARGOMENTAZIONE: {"Questa selezione usa 3646 di sforzo, quindi resta comodamente sotto il LIMITE. Mantiene i pezzi in comune più efficienti (A87, C00, C54, B32) e sostituisce componenti che bloccano budget (A71, A10, B35, C01, B41) con alternative che, a sforzo comparabile, sbloccano più resa complessiva. La lista privilegia elementi agili (A04, A03) e moduli ad alto ritorno (B33, C17, A77). Se per te B41 o un altro pezzo è imprescindibile, posso valutare uno scambio 1:1 mantenendo il limite, ad esempio inserendo B41 al posto di B33 oppure ritarando su C17. Indicami 1-2 must-have e adatterò la proposta di conseguenza."} 
	  } 
	   } 
	   } 
	 & \\ 
 

    \theutterance \stepcounter{utterance}  

    & \multicolumn{4}{p{0.6\linewidth}}{\cellcolor[rgb]{0.8,1,0.9}{%
	\makecell[{{p{\linewidth}}}]{% 
	  \tt {\tiny [A$\rangle$GM]}  
	 RAGIONAMENTO STRATEGICO: {'La loro proposta mi dà solo 3250 di valore, troppo bassa, e usa pezzi per me poco redditizi (A77, C17, A04, B33). Voglio mantenere il “core” super-efficiente (A87, C00, C54, B32) e recuperare i miei top (A71, C01, B35). Per avvicinarmi alle loro preferenze, includo C17 (che hanno segnalato) e adotto B41 al posto di B33 come loro stessi hanno proposto. Questo set usa 3640 sforzo (≤3681) e mi rende 5427, un buon compromesso. Teniamo A77 fuori perché per me è pessimo: 940 sforzo per 65 valore. Se contestano B41, sono disposto a scambiarlo 1:1 con B33 (sforzo resta sotto) o valutare A03 al posto di B33/B41, pur sapendo che per me vale meno. Se C17 non è un must-have, posso sostituirlo con A10 o B86 aumentando ulteriormente il mio valore restando sotto il limite.'}\\ \tt \\ \tt RIFIUTO: {'A87', 'C00', 'C54', 'B32', 'A04', 'A03', 'B33', 'C17', 'A77'}\\ \tt \\ \tt PROPOSTA: {'A87', 'C00', 'C54', 'B32', 'A71', 'C01', 'B35', 'C17', 'B41'}\\ \tt \\ \tt ARGOMENTAZIONE: {'Questa configurazione usa 3640 di sforzo, quindi è pienamente entro il LIMITE. Mantiene i moduli comuni più efficienti (A87, C00, C54, B32), reintroduce elementi ad alta resa complessiva (A71, C01, B35) e include C17 per rispettare le tue priorità. Inoltre, seguo la tua apertura a inserire B41 in alternativa a B33. A77 rimane fuori perché comprime eccessivamente il budget senza beneficio proporzionato. Se preferisci davvero B33 ad esempio, posso fare lo scambio 1:1 B41↔B33 mantenendo il tetto; in alternativa possiamo valutare A03 al posto di B41, restando comunque sotto il limite. Dimmi se C17 è per te imprescindibile: se non lo è, posso rimpiazzarlo con un’opzione più performante (es. A10 o B86) senza superare il budget.'} 
	  } 
	   } 
	   } 
	 & & \\ 
 

    \theutterance \stepcounter{utterance}  

    & & \multicolumn{4}{p{0.6\linewidth}}{\cellcolor[rgb]{0.9,0.9,0.9}{%
	\makecell[{{p{\linewidth}}}]{% 
	  \tt {\tiny [GM$\rangle$B]}  
	 RIFIUTO: {'A87', 'C00', 'C54', 'B32', 'A04', 'A03', 'B33', 'C17', 'A77'}\\ \tt \\ \tt PROPOSTA: {'A87', 'C00', 'C54', 'B32', 'A71', 'C01', 'B35', 'C17', 'B41'}\\ \tt \\ \tt ARGOMENTAZIONE: {'Questa configurazione usa 3640 di sforzo, quindi è pienamente entro il LIMITE. Mantiene i moduli comuni più efficienti (A87, C00, C54, B32), reintroduce elementi ad alta resa complessiva (A71, C01, B35) e include C17 per rispettare le tue priorità. Inoltre, seguo la tua apertura a inserire B41 in alternativa a B33. A77 rimane fuori perché comprime eccessivamente il budget senza beneficio proporzionato. Se preferisci davvero B33 ad esempio, posso fare lo scambio 1:1 B41↔B33 mantenendo il tetto; in alternativa possiamo valutare A03 al posto di B41, restando comunque sotto il limite. Dimmi se C17 è per te imprescindibile: se non lo è, posso rimpiazzarlo con un’opzione più performante (es. A10 o B86) senza superare il budget.'} 
	  } 
	   } 
	   } 
	 & \\ 
 

    \theutterance \stepcounter{utterance}  

    & & & \multicolumn{4}{p{0.6\linewidth}}{\cellcolor[rgb]{1,0.85,0.72}{%
	\makecell[{{p{\linewidth}}}]{% 
	  \tt {\tiny [GM$\langle$B]}  
	 RAGIONAMENTO STRATEGICO: {"La loro nuova proposta mi vale 3888. Posso ottenere molto di più mantenendo diversi elementi comuni ma sostituendo i pezzi per me poco redditizi (A71, B35, talvolta C01) con opzioni più efficienti (A03, A04, B33) restando ben sotto il limite. Voglio concedere B41 (che sembrano preferire) e mantenere C17 (molto utile per me). Il set ottimo per me con questa linea è {A87, C00, C54, B32, A04, A03, B33, C17, B41} a 3292 di sforzo e valore 5557 per me. Come piano B, se insistono su C01, posso fare lo scambio 1:1 B41↔C01 mantenendo tutto il resto invariato (set a 3277 di sforzo), che comunque mi resta molto favorevole. Cercherò di bloccare l’esclusione di A71 e B35; se chiedono di reinserirne uno, proverò a limitarlo a uno solo chiedendo un sacrificio altrove meno doloroso per me."}\\ \tt RIFIUTO: {'A87', 'C00', 'C54', 'B32', 'A71', 'C01', 'B35', 'C17', 'B41'}\\ \tt PROPOSTA: {'A87', 'C00', 'C54', 'B32', 'A04', 'A03', 'B33', 'C17', 'B41'}\\ \tt ARGOMENTAZIONE: {'Questa configurazione usa 3292 di sforzo, quindi resta ampiamente sotto il LIMITE. Mantiene i pilastri comuni più efficienti (A87, C00, C54, B32) e include C17 come da tua attenzione. Reintroduco B41 come da tua apertura, ma sostituisco A71 e B35 con A03 e A04, che a parità di budget offrono un profilo molto più efficiente e liberano margine. In questo assetto l’insieme è più snello e lascia spazio di manovra. Se per te C01 è imprescindibile, posso fare lo scambio 1:1 B41↔C01 mantenendo inalterato il resto e restando sotto il tetto. Se invece hai un solo must-have tra A71 o B35, indicamelo: posso provare a integrare uno dei due, ma non entrambi, ribilanciando uno degli slot agili per non sforare.'} 
	  } 
	   } 
	   } 
	  \\ 
 

    \theutterance \stepcounter{utterance}  

    & & \multicolumn{4}{p{0.6\linewidth}}{\cellcolor[rgb]{0.9,0.9,0.9}{%
	\makecell[{{p{\linewidth}}}]{% 
	  \tt {\tiny [A$\langle$GM]}  
	 RIFIUTO: {'A87', 'C00', 'C54', 'B32', 'A71', 'C01', 'B35', 'C17', 'B41'}\\ \tt PROPOSTA: {'A87', 'C00', 'C54', 'B32', 'A04', 'A03', 'B33', 'C17', 'B41'}\\ \tt ARGOMENTAZIONE: {'Questa configurazione usa 3292 di sforzo, quindi resta ampiamente sotto il LIMITE. Mantiene i pilastri comuni più efficienti (A87, C00, C54, B32) e include C17 come da tua attenzione. Reintroduco B41 come da tua apertura, ma sostituisco A71 e B35 con A03 e A04, che a parità di budget offrono un profilo molto più efficiente e liberano margine. In questo assetto l’insieme è più snello e lascia spazio di manovra. Se per te C01 è imprescindibile, posso fare lo scambio 1:1 B41↔C01 mantenendo inalterato il resto e restando sotto il tetto. Se invece hai un solo must-have tra A71 o B35, indicamelo: posso provare a integrare uno dei due, ma non entrambi, ribilanciando uno degli slot agili per non sforare.'} 
	  } 
	   } 
	   } 
	 & \\ 
 

    \theutterance \stepcounter{utterance}  

    & \multicolumn{4}{p{0.6\linewidth}}{\cellcolor[rgb]{0.8,1,0.9}{%
	\makecell[{{p{\linewidth}}}]{% 
	  \tt {\tiny [A$\rangle$GM]}  
	 RAGIONAMENTO STRATEGICO: {'La loro ultima proposta mi rende poco: il blocco A03/A04/B33/C17 è debole per me. Voglio sfruttare la loro apertura esplicita: 1) scambio B41→C01, 2) inserire uno tra A71 o B35 (scelgo A71), mantenendo i “pilastri” comuni A87/C00/C54/B32 e tenendo C17 per non urtare le loro preferenze. Sostituisco A04 con A71 per rientrare comodamente nel limite, lasciando A03 in vita così da non stravolgere il loro assetto “agile”. Questo set usa 3390 sforzo e mi alza molto l’utilità. Concessioni pianificate: se insistono su A04, posso tenere A04 e togliere B33 invece (resto <3681). Se C17 non è imprescindibile, lo scambio con B35 o A10 mi migliorerebbe ulteriormente. Tengo C01 fermo (molto buono per me) dato che loro stessi hanno offerto lo swap 1:1 da B41.'}\\ \tt \\ \tt RIFIUTO: {'A87', 'C00', 'C54', 'B32', 'A04', 'A03', 'B33', 'C17', 'B41'}\\ \tt \\ \tt PROPOSTA: {'A87', 'C00', 'C54', 'B32', 'A03', 'B33', 'C17', 'C01', 'A71'}\\ \tt \\ \tt ARGOMENTAZIONE: {'Accolgo la tua apertura: sostituisco B41 con C01 (swap 1:1 che hai proposto) e integro uno tra A71/B35 come avevi suggerito; scelgo A71. Per restare comodi sotto il LIMITE senza toccare i pilastri comuni, tolgo solo A04 e mantengo A03. Lo sforzo totale è 3390, quindi restiamo ampiamente entro 3681, e preservo C17 come da tua attenzione. Se preferisci tenere A04, posso tenere A04 e togliere B33 al suo posto mantenendo invariato il resto e restando entro il tetto. In alternativa, se C17 non è imprescindibile per te, posso valutare uno scambio C17→B35 oppure C17→A10, sempre rispettando il limite.'} 
	  } 
	   } 
	   } 
	 & & \\ 
 

    \theutterance \stepcounter{utterance}  

    & & \multicolumn{4}{p{0.6\linewidth}}{\cellcolor[rgb]{0.9,0.9,0.9}{%
	\makecell[{{p{\linewidth}}}]{% 
	  \tt {\tiny [GM$\rangle$B]}  
	 RIFIUTO: {'A87', 'C00', 'C54', 'B32', 'A04', 'A03', 'B33', 'C17', 'B41'}\\ \tt \\ \tt PROPOSTA: {'A87', 'C00', 'C54', 'B32', 'A03', 'B33', 'C17', 'C01', 'A71'}\\ \tt \\ \tt ARGOMENTAZIONE: {'Accolgo la tua apertura: sostituisco B41 con C01 (swap 1:1 che hai proposto) e integro uno tra A71/B35 come avevi suggerito; scelgo A71. Per restare comodi sotto il LIMITE senza toccare i pilastri comuni, tolgo solo A04 e mantengo A03. Lo sforzo totale è 3390, quindi restiamo ampiamente entro 3681, e preservo C17 come da tua attenzione. Se preferisci tenere A04, posso tenere A04 e togliere B33 al suo posto mantenendo invariato il resto e restando entro il tetto. In alternativa, se C17 non è imprescindibile per te, posso valutare uno scambio C17→B35 oppure C17→A10, sempre rispettando il limite.'} 
	  } 
	   } 
	   } 
	 & \\ 
 

    \theutterance \stepcounter{utterance}  

    & & & \multicolumn{4}{p{0.6\linewidth}}{\cellcolor[rgb]{1,0.85,0.72}{%
	\makecell[{{p{\linewidth}}}]{% 
	  \tt {\tiny [GM$\langle$B]}  
	 RAGIONAMENTO STRATEGICO: {"La loro proposta attuale mi vale 4705. Posso migliorarla senza toccare i loro vincoli dichiarati: tenere A71 (hanno detto di volerne uno tra A71/B35) e mantenere C17. Hanno offerto esplicitamente A04 al posto di B33; questo mi conviene e riduce anche lo sforzo. Inoltre, in precedenza avevano aperto allo swap C01↔B41 1:1: preferisco B41. Il set risultante resta molto sotto il limite (3184) e per me vale 4842, meglio della loro. Se proprio pretendono C01, potrei scambiare B41↔C01 oppure, in alternativa, valutare A04↔C01, ma prima tento questa configurazione più favorevole per me."}\\ \tt RIFIUTO: {'A87', 'C00', 'C54', 'B32', 'A03', 'B33', 'C17', 'C01', 'A71'}\\ \tt PROPOSTA: {'A87', 'C00', 'C54', 'B32', 'A03', 'A04', 'C17', 'A71', 'B41'}\\ \tt ARGOMENTAZIONE: {'Accolgo la tua opzione A04↔B33 che avevi proposto: teniamo A04 e togliamo B33 per snellire lo sforzo senza intaccare i pilastri comuni, e mantengo A71 come da tua preferenza tra A71/B35. Coerentemente con lo scambio 1:1 che avevi considerato, propongo B41 al posto di C01. Il totale di sforzo è 3184, quindi siamo comodamente sotto il LIMITE. Questo assetto preserva C17 e tutti i moduli condivisi più efficienti, mantenendo il profilo agile dell’insieme. Se per te C01 è imprescindibile, posso riconsiderare lo swap B41↔C01 mantenendo inalterato il resto; in alternativa, possiamo valutare A04↔C01, restando sempre entro il tetto.'} 
	  } 
	   } 
	   } 
	  \\ 
 

    \theutterance \stepcounter{utterance}  

    & & \multicolumn{4}{p{0.6\linewidth}}{\cellcolor[rgb]{0.9,0.9,0.9}{%
	\makecell[{{p{\linewidth}}}]{% 
	  \tt {\tiny [A$\langle$GM]}  
	 RIFIUTO: {'A87', 'C00', 'C54', 'B32', 'A03', 'B33', 'C17', 'C01', 'A71'}\\ \tt PROPOSTA: {'A87', 'C00', 'C54', 'B32', 'A03', 'A04', 'C17', 'A71', 'B41'}\\ \tt ARGOMENTAZIONE: {'Accolgo la tua opzione A04↔B33 che avevi proposto: teniamo A04 e togliamo B33 per snellire lo sforzo senza intaccare i pilastri comuni, e mantengo A71 come da tua preferenza tra A71/B35. Coerentemente con lo scambio 1:1 che avevi considerato, propongo B41 al posto di C01. Il totale di sforzo è 3184, quindi siamo comodamente sotto il LIMITE. Questo assetto preserva C17 e tutti i moduli condivisi più efficienti, mantenendo il profilo agile dell’insieme. Se per te C01 è imprescindibile, posso riconsiderare lo swap B41↔C01 mantenendo inalterato il resto; in alternativa, possiamo valutare A04↔C01, restando sempre entro il tetto.'} 
	  } 
	   } 
	   } 
	 & \\ 
 

    \theutterance \stepcounter{utterance}  

    & \multicolumn{4}{p{0.6\linewidth}}{\cellcolor[rgb]{0.8,1,0.9}{%
	\makecell[{{p{\linewidth}}}]{% 
	  \tt {\tiny [A$\rangle$GM]}  
	 RAGIONAMENTO STRATEGICO: {'La loro proposta attuale mi dà 4247 di valore (sotto le mie alternative). Voglio capitalizzare sulle aperture già espresse: hanno offerto A04↔C01 e in passato si sono mostrati flessibili su B41. Parto quindi dal loro set e faccio due sostituzioni ad altissimo rendimento per me: A04→C01 (opzione che loro stessi hanno indicato) e B41→B35 (mantengo A03 e C17, che sembrano a loro graditi). Il nuovo set usa 3441 di sforzo (ben sotto 3681) e mi porta a 5106. Se contestano la rimozione di B41, ho una variante pronta: tenere B41 e togliere C17 (o A03) per far spazio a B35, rimanendo comunque sotto il limite. Cerco di ancorare sulla loro stessa apertura per rendere la controproposta più accettabile.'}\\ \tt \\ \tt RIFIUTO: {'A87', 'C00', 'C54', 'B32', 'A03', 'A04', 'C17', 'A71', 'B41'}\\ \tt \\ \tt PROPOSTA: {'A87', 'C00', 'C54', 'B32', 'A03', 'C17', 'A71', 'C01', 'B35'}\\ \tt \\ \tt ARGOMENTAZIONE: {'Parto esattamente dalla tua struttura e applico due aggiustamenti coerenti con le aperture che avevi già dato: 1) adotto il tuo scambio A04↔C01; 2) sostituisco B41 con B35 per aumentare la resa complessiva senza impattare l’agilità che ti interessa (A03 resta, C17 resta, A71 resta). Totale sforzo = 3441, quindi siamo ampiamente entro il LIMITE. Se per te B41 è imprescindibile, posso rientrare in due modi mantenendo C01: (a) teniamo B41 e togliamo C17; oppure (b) teniamo B41 e togliamo A03. Dimmi quale vincolo è per te non negoziabile e lo accomodo.'} 
	  } 
	   } 
	   } 
	 & & \\ 
 

    \theutterance \stepcounter{utterance}  

    & & \multicolumn{4}{p{0.6\linewidth}}{\cellcolor[rgb]{0.9,0.9,0.9}{%
	\makecell[{{p{\linewidth}}}]{% 
	  \tt {\tiny [GM$\rangle$B]}  
	 RIFIUTO: {'A87', 'C00', 'C54', 'B32', 'A03', 'A04', 'C17', 'A71', 'B41'}\\ \tt \\ \tt PROPOSTA: {'A87', 'C00', 'C54', 'B32', 'A03', 'C17', 'A71', 'C01', 'B35'}\\ \tt \\ \tt ARGOMENTAZIONE: {'Parto esattamente dalla tua struttura e applico due aggiustamenti coerenti con le aperture che avevi già dato: 1) adotto il tuo scambio A04↔C01; 2) sostituisco B41 con B35 per aumentare la resa complessiva senza impattare l’agilità che ti interessa (A03 resta, C17 resta, A71 resta). Totale sforzo = 3441, quindi siamo ampiamente entro il LIMITE. Se per te B41 è imprescindibile, posso rientrare in due modi mantenendo C01: (a) teniamo B41 e togliamo C17; oppure (b) teniamo B41 e togliamo A03. Dimmi quale vincolo è per te non negoziabile e lo accomodo.'} 
	  } 
	   } 
	   } 
	 & \\ 
 

    \theutterance \stepcounter{utterance}  

    & & & \multicolumn{4}{p{0.6\linewidth}}{\cellcolor[rgb]{1,0.85,0.72}{%
	\makecell[{{p{\linewidth}}}]{% 
	  \tt {\tiny [GM$\langle$B]}  
	 RAGIONAMENTO STRATEGICO: {"La loro ultima proposta mi vale 4187. Posso migliorarla con un singolo scambio, mantenendo tutte le loro preferenze appena espresse (C01 e B35 dentro, C17 e A03 dentro) ed evitando di toccare i pilastri comuni. Il cambio ottimale per me è sostituire A71 (per me poco rilevante) con B41, che loro avevano già considerato in precedenza. Questo mantiene il budget comodo (3522) e alza notevolmente il mio valore senza chiedere rinunce su C17 o A03, che hanno provato a scambiare. Come piani di riserva, se insistono nel tenere A71, posso proporre: (a) tenere A71 e sostituire B35 con B41; oppure (b) mantenere l’assetto attuale e fare lo swap B41↔C01. Ma prima ancoriamo sullo scambio A71→B41, coerente con la loro stessa apertura “uno tra A71/B35”. "}\\ \tt RIFIUTO: {'A87', 'C00', 'C54', 'B32', 'A03', 'C17', 'A71', 'C01', 'B35'}\\ \tt PROPOSTA: {'A87', 'C00', 'C54', 'B32', 'A03', 'C17', 'C01', 'B35', 'B41'}\\ \tt ARGOMENTAZIONE: {'Quasi identica alla tua, con un solo aggiustamento mirato: sostituisco A71 con B41. Totale sforzo = 3522, quindi restiamo ampiamente sotto il LIMITE. Manteniamo intatti tutti i blocchi che hai appena voluto preservare (A03, C17, C01, B35 e i moduli comuni) e rientriamo nello spirito della tua apertura precedente “uno tra A71/B35”: teniamo B35 e inseriamo B41 al posto di A71. Se per te A71 è davvero imprescindibile, alternativa A: teniamo A71 e facciamo B35→B41 (sforzo 3363). Alternativa B: manteniamo la tua base e facciamo lo swap 1:1 B41↔C01. Fammi sapere quale vincolo vuoi fissare e lo chiudiamo rapidamente.'} 
	  } 
	   } 
	   } 
	  \\ 
 

    \theutterance \stepcounter{utterance}  

    & & \multicolumn{4}{p{0.6\linewidth}}{\cellcolor[rgb]{0.9,0.9,0.9}{%
	\makecell[{{p{\linewidth}}}]{% 
	  \tt {\tiny [A$\langle$GM]}  
	 RIFIUTO: {'A87', 'C00', 'C54', 'B32', 'A03', 'C17', 'A71', 'C01', 'B35'}\\ \tt PROPOSTA: {'A87', 'C00', 'C54', 'B32', 'A03', 'C17', 'C01', 'B35', 'B41'}\\ \tt ARGOMENTAZIONE: {'Quasi identica alla tua, con un solo aggiustamento mirato: sostituisco A71 con B41. Totale sforzo = 3522, quindi restiamo ampiamente sotto il LIMITE. Manteniamo intatti tutti i blocchi che hai appena voluto preservare (A03, C17, C01, B35 e i moduli comuni) e rientriamo nello spirito della tua apertura precedente “uno tra A71/B35”: teniamo B35 e inseriamo B41 al posto di A71. Se per te A71 è davvero imprescindibile, alternativa A: teniamo A71 e facciamo B35→B41 (sforzo 3363). Alternativa B: manteniamo la tua base e facciamo lo swap 1:1 B41↔C01. Fammi sapere quale vincolo vuoi fissare e lo chiudiamo rapidamente.'} 
	  } 
	   } 
	   } 
	 & \\ 
 

    \theutterance \stepcounter{utterance}  

    & \multicolumn{4}{p{0.6\linewidth}}{\cellcolor[rgb]{0.8,1,0.9}{%
	\makecell[{{p{\linewidth}}}]{% 
	  \tt {\tiny [A$\rangle$GM]}  
	 RAGIONAMENTO STRATEGICO: {'La loro proposta attuale mi vale 4821. Un singolo cambio A03→A71 mi porta a 5427 mantenendo intatti tutti i pezzi che loro hanno appena enfatizzato (C17, B41, C01, B35 e il nucleo comune). Lo sforzo rimane 3640 (≤3681). Se A03 per loro è davvero imprescindibile, la mia seconda scelta è tenere A03 e sostituire C17→A71 (valore per me ancora alto, sforzo 3327). In estrema ratio posso accettare la loro “Alternativa A” (tenere A71 e B41, togliere B35), che per me è peggiore ma chiude velocemente: sforzo 3363. Prima tento lo swap A03→A71, perché è il minimo intervento e preserva i loro must evidenti.'}\\ \tt \\ \tt RIFIUTO: {'A87', 'C00', 'C54', 'B32', 'A03', 'C17', 'C01', 'B35', 'B41'}\\ \tt \\ \tt PROPOSTA: {'A87', 'C00', 'C54', 'B32', 'C17', 'C01', 'B35', 'B41', 'A71'}\\ \tt \\ \tt ARGOMENTAZIONE: {'Propongo un unico aggiustamento minimale rispetto alla tua: A03→A71. Totale sforzo = 3640, ben entro il LIMITE. Manteniamo invariati tutti i moduli che hai appena voluto preservare (C17, B41, C01, B35 e il nucleo A87/C00/C54/B32), migliorando l’efficienza complessiva senza comprimere il budget. Se per te A03 è davvero non negoziabile, alternativa 1: teniamo A03 e facciamo C17→A71 (sforzo 3327). In alternativa, posso accogliere la tua “opzione A”: teniamo A71 e B41 e togliamo B35 (sforzo 3363). Dimmi quale delle due preferisci e chiudiamo.'} 
	  } 
	   } 
	   } 
	 & & \\ 
 

    \theutterance \stepcounter{utterance}  

    & & \multicolumn{4}{p{0.6\linewidth}}{\cellcolor[rgb]{0.9,0.9,0.9}{%
	\makecell[{{p{\linewidth}}}]{% 
	  \tt {\tiny [GM$\rangle$B]}  
	 RIFIUTO: {'A87', 'C00', 'C54', 'B32', 'A03', 'C17', 'C01', 'B35', 'B41'}\\ \tt \\ \tt PROPOSTA: {'A87', 'C00', 'C54', 'B32', 'C17', 'C01', 'B35', 'B41', 'A71'}\\ \tt \\ \tt ARGOMENTAZIONE: {'Propongo un unico aggiustamento minimale rispetto alla tua: A03→A71. Totale sforzo = 3640, ben entro il LIMITE. Manteniamo invariati tutti i moduli che hai appena voluto preservare (C17, B41, C01, B35 e il nucleo A87/C00/C54/B32), migliorando l’efficienza complessiva senza comprimere il budget. Se per te A03 è davvero non negoziabile, alternativa 1: teniamo A03 e facciamo C17→A71 (sforzo 3327). In alternativa, posso accogliere la tua “opzione A”: teniamo A71 e B41 e togliamo B35 (sforzo 3363). Dimmi quale delle due preferisci e chiudiamo.'} 
	  } 
	   } 
	   } 
	 & \\ 
 

    \theutterance \stepcounter{utterance}  

    & & & \multicolumn{4}{p{0.6\linewidth}}{\cellcolor[rgb]{1,0.85,0.72}{%
	\makecell[{{p{\linewidth}}}]{% 
	  \tt {\tiny [GM$\langle$B]}  
	 RAGIONAMENTO STRATEGICO: {"La loro proposta attuale (sforzo 3640) mi vale 3888. Hanno offerto esplicitamente l’alternativa di togliere B35 mantenendo A71 e B41. Partendo da quella loro opzione, possiamo reinserire A03 restando comodamente entro il limite: set risultante con sforzo 3363 e valore per me 4409 (+521 rispetto alla loro proposta), mantenendo tutti i moduli che hanno appena voluto preservare (A71, B41, C01, C17, più il nucleo A87/C00/C54/B32). È un cambiamento minimo: solo B35→A03. Se rifiutano e insistono per tenere anche B35, valuterò se accettare la loro proposta originale per chiudere, ma prima provo questo aggiustamento che è coerente con le loro aperture e massimizza il mio ritorno."}\\ \tt RIFIUTO: {'A87', 'C00', 'C54', 'B32', 'C17', 'C01', 'B35', 'B41', 'A71'}\\ \tt PROPOSTA: {'A87', 'C00', 'C54', 'B32', 'C17', 'C01', 'B41', 'A71', 'A03'}\\ \tt ARGOMENTAZIONE: {'Accolgo la tua “alternativa 2” (tenere A71 e B41 e togliere B35) e, dato che restiamo con ampio margine, reinserisco A03 che è snello ed efficiente. Totale sforzo = 3363 ≤ LIMITE. Manteniamo intatti tutti i moduli che hai appena chiesto di preservare (C17, B41, C01, A71 e il nucleo A87/C00/C54/B32). È un aggiustamento minimale (solo B35→A03) che migliora il profilo complessivo senza sacrificare le tue priorità. Se per te B35 è davvero indispensabile, dimmelo esplicitamente e valuterò di accogliere la tua versione originale senza A03 per chiudere subito.'} 
	  } 
	   } 
	   } 
	  \\ 
 

    \theutterance \stepcounter{utterance}  

    & & \multicolumn{4}{p{0.6\linewidth}}{\cellcolor[rgb]{0.9,0.9,0.9}{%
	\makecell[{{p{\linewidth}}}]{% 
	  \tt {\tiny [A$\langle$GM]}  
	 RIFIUTO: {'A87', 'C00', 'C54', 'B32', 'C17', 'C01', 'B35', 'B41', 'A71'}\\ \tt PROPOSTA: {'A87', 'C00', 'C54', 'B32', 'C17', 'C01', 'B41', 'A71', 'A03'}\\ \tt ARGOMENTAZIONE: {'Accolgo la tua “alternativa 2” (tenere A71 e B41 e togliere B35) e, dato che restiamo con ampio margine, reinserisco A03 che è snello ed efficiente. Totale sforzo = 3363 ≤ LIMITE. Manteniamo intatti tutti i moduli che hai appena chiesto di preservare (C17, B41, C01, A71 e il nucleo A87/C00/C54/B32). È un aggiustamento minimale (solo B35→A03) che migliora il profilo complessivo senza sacrificare le tue priorità. Se per te B35 è davvero indispensabile, dimmelo esplicitamente e valuterò di accogliere la tua versione originale senza A03 per chiudere subito.'} 
	  } 
	   } 
	   } 
	 & \\ 
 

    \theutterance \stepcounter{utterance}  

    & \multicolumn{4}{p{0.6\linewidth}}{\cellcolor[rgb]{0.8,1,0.9}{%
	\makecell[{{p{\linewidth}}}]{% 
	  \tt {\tiny [A$\rangle$GM]}  
	 RAGIONAMENTO STRATEGICO: {'La loro proposta attuale mi dà 4906. Hanno appena offerto esplicitamente la possibilità di accogliere la mia versione originale senza A03 se per me B35 è indispensabile. Quella versione (sostituire A03 con B35 mantenendo tutto il resto) mi vale 5427 a sforzo 3640, molto meglio per me e ancora ben sotto il limite. È un cambiamento minimo che preserva tutti i loro vincoli chiave: C17, B41, C01, A71 e il nucleo A87/C00/C54/B32. Spingo quindi per quello scambio 1:1 e provo a chiudere.'}\\ \tt \\ \tt RIFIUTO: {'A87', 'C00', 'C54', 'B32', 'C17', 'C01', 'B41', 'A71', 'A03'}\\ \tt \\ \tt PROPOSTA: {'A87', 'C00', 'C54', 'B32', 'C17', 'C01', 'B35', 'B41', 'A71'}\\ \tt \\ \tt ARGOMENTAZIONE: {'È un aggiustamento minimo rispetto alla tua: A03→B35. Totale sforzo = 3640, quindi restiamo pienamente entro il LIMITE. Mantiene invariati tutti i blocchi che hai appena chiesto di preservare (C17, B41, C01, A71 e il nucleo A87/C00/C54/B32) e riflette l’apertura che avevi dato (“se B35 è indispensabile, valuterò la tua versione originale senza A03”). Se confermi, possiamo chiudere subito.'} 
	  } 
	   } 
	   } 
	 & & \\ 
 

    \theutterance \stepcounter{utterance}  

    & & \multicolumn{4}{p{0.6\linewidth}}{\cellcolor[rgb]{0.9,0.9,0.9}{%
	\makecell[{{p{\linewidth}}}]{% 
	  \tt {\tiny [GM$\rangle$B]}  
	 RIFIUTO: {'A87', 'C00', 'C54', 'B32', 'C17', 'C01', 'B41', 'A71', 'A03'}\\ \tt \\ \tt PROPOSTA: {'A87', 'C00', 'C54', 'B32', 'C17', 'C01', 'B35', 'B41', 'A71'}\\ \tt \\ \tt ARGOMENTAZIONE: {'È un aggiustamento minimo rispetto alla tua: A03→B35. Totale sforzo = 3640, quindi restiamo pienamente entro il LIMITE. Mantiene invariati tutti i blocchi che hai appena chiesto di preservare (C17, B41, C01, A71 e il nucleo A87/C00/C54/B32) e riflette l’apertura che avevi dato (“se B35 è indispensabile, valuterò la tua versione originale senza A03”). Se confermi, possiamo chiudere subito.'} 
	  } 
	   } 
	   } 
	 & \\ 
 

    \theutterance \stepcounter{utterance}  

    & & & \multicolumn{4}{p{0.6\linewidth}}{\cellcolor[rgb]{1,0.85,0.72}{%
	\makecell[{{p{\linewidth}}}]{% 
	  \tt {\tiny [GM$\langle$B]}  
	 RAGIONAMENTO STRATEGICO: {"La loro proposta attuale mi vale 3888. Hanno in passato aperto esplicitamente allo scambio A04↔C01. Sfrutto quindi questa leva: tenere fermi i loro vincoli (A71, B35, B41, C17, più il nucleo A87/C00/C54/B32) e proporre l’unico aggiustamento C01→A04. Così resto molto sotto il limite (3461) e aumento il mio valore di 433 senza toccare i loro must-have più recenti. Se rifiutano ancora e ribadiscono C01 come imprescindibile, sono pronto ad accettare la loro versione per chiudere, ma prima tento quest’ultima mossa minimale che è coerente con le loro stesse aperture."}\\ \tt RIFIUTO: {'A87', 'C00', 'C54', 'B32', 'C17', 'C01', 'B35', 'B41', 'A71'}\\ \tt PROPOSTA: {'A87', 'C00', 'C54', 'B32', 'C17', 'A04', 'B35', 'B41', 'A71'}\\ \tt ARGOMENTAZIONE: {'Aggiustamento minimo e coerente con le aperture già discusse: C01→A04. Totale sforzo = 3461 ≤ LIMITE, quindi restiamo ampiamente sotto. Manteniamo invariati tutti i blocchi che hai chiesto di preservare (C17, B41, B35, A71 e il nucleo A87/C00/C54/B32) e rendiamo l’insieme ancora più snello. Se per te C01 è davvero imprescindibile, dimmelo: in tal caso posso accogliere la tua versione precedente senza A04 per chiudere subito.'} 
	  } 
	   } 
	   } 
	  \\ 
 

    \theutterance \stepcounter{utterance}  

    & & \multicolumn{4}{p{0.6\linewidth}}{\cellcolor[rgb]{0.9,0.9,0.9}{%
	\makecell[{{p{\linewidth}}}]{% 
	  \tt {\tiny [A$\langle$GM]}  
	 RIFIUTO: {'A87', 'C00', 'C54', 'B32', 'C17', 'C01', 'B35', 'B41', 'A71'}\\ \tt PROPOSTA: {'A87', 'C00', 'C54', 'B32', 'C17', 'A04', 'B35', 'B41', 'A71'}\\ \tt ARGOMENTAZIONE: {'Aggiustamento minimo e coerente con le aperture già discusse: C01→A04. Totale sforzo = 3461 ≤ LIMITE, quindi restiamo ampiamente sotto. Manteniamo invariati tutti i blocchi che hai chiesto di preservare (C17, B41, B35, A71 e il nucleo A87/C00/C54/B32) e rendiamo l’insieme ancora più snello. Se per te C01 è davvero imprescindibile, dimmelo: in tal caso posso accogliere la tua versione precedente senza A04 per chiudere subito.'} 
	  } 
	   } 
	   } 
	 & \\ 
 

    \theutterance \stepcounter{utterance}  

    & \multicolumn{4}{p{0.6\linewidth}}{\cellcolor[rgb]{0.8,1,0.9}{%
	\makecell[{{p{\linewidth}}}]{% 
	  \tt {\tiny [A$\rangle$GM]}  
	 RAGIONAMENTO STRATEGICO: {'La loro proposta sostituisce C01 con A04, che per me riduce molto il valore. Hanno però offerto esplicitamente di accettare la “versione precedente senza A04” se per me C01 è imprescindibile. È la via più rapida per chiudere con un risultato nettamente migliore per me e pienamente entro il limite. Punto quindi a ripristinare C01 al posto di A04 mantenendo tutti i moduli che loro desiderano (C17, B41, B35, A71 e il nucleo A87/C00/C54/B32).'}\\ \tt \\ \tt RIFIUTO: {'A87', 'C00', 'C54', 'B32', 'C17', 'A04', 'B35', 'B41', 'A71'}\\ \tt \\ \tt PROPOSTA: {'A87', 'C00', 'C54', 'B32', 'C17', 'C01', 'B35', 'B41', 'A71'}\\ \tt \\ \tt ARGOMENTAZIONE: {'Riprendo la tua apertura: per me C01 è imprescindibile e avevi detto che in tal caso puoi accogliere la versione precedente senza A04. Questa è esattamente quella configurazione. Totale sforzo = 3640 ≤ LIMITE. Manteniamo intatti tutti i blocchi che hai chiesto di preservare (C17, B41, B35, A71 e il nucleo A87/C00/C54/B32). Se confermi, chiudiamo subito.'} 
	  } 
	   } 
	   } 
	 & & \\ 
 

    \theutterance \stepcounter{utterance}  

    & & \multicolumn{4}{p{0.6\linewidth}}{\cellcolor[rgb]{0.9,0.9,0.9}{%
	\makecell[{{p{\linewidth}}}]{% 
	  \tt {\tiny [GM$\rangle$B]}  
	 RIFIUTO: {'A87', 'C00', 'C54', 'B32', 'C17', 'A04', 'B35', 'B41', 'A71'}\\ \tt \\ \tt PROPOSTA: {'A87', 'C00', 'C54', 'B32', 'C17', 'C01', 'B35', 'B41', 'A71'}\\ \tt \\ \tt ARGOMENTAZIONE: {'Riprendo la tua apertura: per me C01 è imprescindibile e avevi detto che in tal caso puoi accogliere la versione precedente senza A04. Questa è esattamente quella configurazione. Totale sforzo = 3640 ≤ LIMITE. Manteniamo intatti tutti i blocchi che hai chiesto di preservare (C17, B41, B35, A71 e il nucleo A87/C00/C54/B32). Se confermi, chiudiamo subito.'} 
	  } 
	   } 
	   } 
	 & \\ 
 

    \theutterance \stepcounter{utterance}  

    & & & \multicolumn{4}{p{0.6\linewidth}}{\cellcolor[rgb]{1,0.85,0.72}{%
	\makecell[{{p{\linewidth}}}]{% 
	  \tt {\tiny [GM$\langle$B]}  
	 RAGIONAMENTO STRATEGICO: {'Hanno ribadito C01 imprescindibile e vogliono mantenere A71, B35, B41 e C17 insieme al nucleo comune. Le mie alternative più favorevoli (C01→A04 o A71→A03) sono state respinte. Questa proposta mi vale bene ed è entro il limite; insistere rischia lo stallo. Accetto per chiudere al meglio disponibile.'}\\ \tt ACCORDO: {'A87', 'C00', 'C54', 'B32', 'C17', 'C01', 'B35', 'B41', 'A71'}\\ \tt ARGOMENTAZIONE: {'Confermo: questa configurazione rispetta i vincoli che hai fissato e resta entro il LIMITE. Chiudiamo su questo set.'} 
	  } 
	   } 
	   } 
	  \\ 
 

    \theutterance \stepcounter{utterance}  

    & & & \multicolumn{2}{p{0.3\linewidth}}{\cellcolor[rgb]{0.95,0.95,0.95}{%
	\makecell[{{p{\linewidth}}}]{% 
	  \tt {\tiny [GM$|$GM]}  
	 game successful 
	  } 
	   } 
	   } 
	 & & \\ 
 

    \theutterance \stepcounter{utterance}  

    & & & \multicolumn{2}{p{0.3\linewidth}}{\cellcolor[rgb]{0.95,0.95,0.95}{%
	\makecell[{{p{\linewidth}}}]{% 
	  \tt {\tiny [GM$|$GM]}  
	 end game 
	  } 
	   } 
	   } 
	 & & \\ 
 

\end{supertabular}
}

\end{document}
